\documentclass[letter,11pt]{article}
\usepackage{hyperref}
\usepackage{graphicx}
\usepackage{fancyhdr}
\usepackage{multicol}
\pagestyle{fancy}
\usepackage[letterpaper, margin=1in]{geometry}
\geometry{letterpaper}
\usepackage{parskip} % Disable initial indent

\usepackage[utf8]{inputenc}
\fancyhf{}
\renewcommand{\headrulewidth}{0pt} % Remove default underline from header package
\rhead{CYBR/CMSC 691 - Lab 12}
%\rhead{}
\lhead{\begin{picture}(0,0) \put(0,-10){\includegraphics[width=1.1cm]{Images/UMBC-vertical}} \end{picture}}
\cfoot{\thepage}
\rfoot{Spring 2020}
\lfoot{Lab 12}
%\AtEndDocument{\vfill \hfill \LaTeX}

\begin{document}

\section*{Lab 12 - PDF/RTF Malware Models}
\paragraph{}Building from Lab 10 on $n$-grams, create a model which detects PE32, and PDF or RTF malware. This would result in a separate model for each file type.

\textbf{Points:} 20

\section*{Instructions}
\begin{enumerate}
    \item Write a script to break the data into separate training and testing directories, with an 80\%/20\% split. You should use symbolic links instead of copying the files, since it will save disk space. In Python: \texttt{os.symlink(existingFilePath, newLinkPath)}.
    \item Make a dataset in CSV format.
    \item Use \texttt{Pandas} to read the CSV file as a dataframe.
    \item Train the model.
    \item Evaluate on the test data.
\end{enumerate}

\section*{Questions}
\paragraph{}In a separate document, provide your answers to the following questions. The expected format is a Word document, or similar, with a few images demonstrating your progress. The assignment will be submitted via Blackboard.
\begin{enumerate}
    \item What was the highest accuracy score you were able to achieve?
    \item Were you able to make the model(s) without code changes?
    \item How many, and what value of $n$ did you use for your $n$-grams?
    \item What troubles or difficulties did you encounter, and how did you resolve them?
\end{enumerate}


\end{document}